\subsection{Электромагнитная индукция}

%%%%%%%%%%%%%%%%%%%%%%%%%%%%%%%%%%%%%%%%%%%%%%%%%%%%
%%%%%%%%%%%%%%%%%%%%%%%%%%%%%%%%%%%%%%%%%%%%%%%%%%%%
 \z{Поперёк магнитного поля индукции $0.1\,\text{Тл}$ движется со скоростью $1\,\text{м/с}$ прямой провод длины $0.3\,\text{м}$. Чему равно напряжение электрического поля между концами проводника?}
%
{$u=Bvl=0.03\,\text{В}$}

%%%%%%%%%%%%%%%%%%%%%%%%%%%%%%%%%%%%%%%%%%%%%%%%%%%%
 \z{При включении магнитного поля, перпендикулярного плоскости витка радиусом $R$, по витку протёк заряд $Q$. Какой заряд протечёт по витку, если его (при неизменном поле) сложить восьмёркой, состоящей из двух окружностей, причём радиус меньшей окружности равен $\frac{R}{4}$? Плоскость восьмёрки также перпендикулярна магнитному полю.}
%
{$-\frac{3}{8}Q;~\frac{1}{2}Q;~\frac{3}{2}Q;~\frac{13}{8}Q$}

%%%%%%%%%%%%%%%%%%%%%%%%%%%%%%%%%%%%%%%%%%%%%%%%%%%%
 \z[{\begin{wrapfigure}[6]{r}{0.2\textwidth}%
		\vspace{-15pt}%
		\begin{flushright}%
 			\includegraphics[width=0.2\textwidth]{5.4.3}%
		\end{flushright}%
	\end{wrapfigure}%
	}]%
{В однородном магнитном поле с индукцией $B=10^{-2}\,\text{Тл}$ расположены вертикально на расстоянии $L=50\,\text{см}$ два металлических прута, замкнутых наверху. По прутьям без трения и без нарушения контакта скользит вниз с постоянной скоростью $v=1\,\text{м/с}$ перемычка $ab$ массой $m=1\,\text{г}$. Определите сопротивление перемычки. Сопротивлением остальной части системы пренебречь.}
%
{$R=\frac{B^2L^2v}{mg} = 2.5 \cdot 10^{-3} \,\text{Ом}$}

%%%%%%%%%%%%%%%%%%%%%%%%%%%%%%%%%%%%%%%%%%%%%%%%%%%%
 \z[{\begin{wrapfigure}[5]{r}{0.2\textwidth}%
		\vspace{-25pt}%
		\begin{flushright}%
 			\includegraphics[width=0.2\textwidth]{5.4.4}%
		\end{flushright}%
	\end{wrapfigure}%
	}]%
{Проволочная квадратная рамка со стороной $l$ свободно падает в поле силы тяжести и магнитном поле, перпендикулярном плоскости рамки, меняющемся по закону $B(z)=B_0+k z$, где $k$ – постоянный коэффициент. Найдите установившуюся скорость рамки.}
%
{$v=\frac{mgR}{k^2l^4}$}

%%%%%%%%%%%%%%%%%%%%%%%%%%%%%%%%%%%%%%%%%%%%%%%%%%%%
 \z[{\begin{wrapfigure}[5]{r}{0.2\textwidth}%
		\vspace{-25pt}%
		\begin{flushright}%
 			\includegraphics[width=0.2\textwidth]{5.4.5}%
		\end{flushright}%
	\end{wrapfigure}%
	}]%
{Прямоугольная проволочная рамка со стороной $L$ находится в магнитном поле с индукцией $B$, перпендикулярной плоскости рамки. По рамке параллельно одной из её сторон без нарушения контакта скользит с постоянной скоростью $v$ перемычка $ab$, сопротивление которой равно $R$. Определите силу тока через перемычку. Сопротивлением остальной части рамки пренебречь.}
%
{$I=\frac{Bvl}{R}$}

%%%%%%%%%%%%%%%%%%%%%%%%%%%%%%%%%%%%%%%%%%%%%%%%%%%%
 \z{Параллельные рельсы находятся в магнитном поле, направление вектора магнитной индукции перпендикулярно плоскости рельсов, индукция поля $B_0$. Рельсы замкнуты перемычкой, имеющей сопротивление $R$. Ещё одна перемычка, имеющая совсем малое сопротивление, может двигаться вдоль рельсов поступательно, замыкая рельсы между собой. С какой силой нужно действовать на подвижную перемычку, чтобы её скорость всё время была $V_0$? Расстояние между рельсами $d$.}
%
{$F=\frac{B^2l^2v_0}{R}$}

%%%%%%%%%%%%%%%%%%%%%%%%%%%%%%%%%%%%%%%%%%%%%%%%%%%%
 \z{Контур из куска тонкого провода сопротивлением $200\,\Omega$ сделан в виде квадрата площади $0.5\,\text{м}^2$. Он находится в магнитном поле $0.2\,\text{Тл}$, вектор магнитной индукции перпендикулярен плоскости контура. Контур деформируют, сделав из него окружность, лежащую в той же плоскости. Найти протекший по контуру заряд.}
%
{$Q=\frac{|\Delta\Phi|}{R}=\frac{B S_0}{R}\left(\frac{4}{\pi}-1\right) \approx 1.37 \cdot 10^{-4}\,\text{Кл}$}




